% ===== PRÉAMBULE CANONIQUE — à coller VERBATIM en tête de CHAQUE .tex =====
\documentclass[11pt,a4paper]{article}
\usepackage[utf8]{inputenc}\usepackage[T1]{fontenc}\usepackage{lmodern}
\usepackage[french]{babel}
\usepackage{amsmath,amssymb,mathtools}\usepackage{icomma}
\usepackage{siunitx}\sisetup{locale=FR}  % \qty{}{}, \si{}, \ang{}
\usepackage{xcolor}\usepackage{colortbl}
\usepackage{tikz}\usepackage{pgfplots}\pgfplotsset{compat=1.18}
\usepackage[siunitx,europeanresistors]{circuitikz} % circuits (élec.)
\usepackage[most]{tcolorbox}\usepackage{enumitem}\usepackage{array,booktabs}
\usepackage[a4paper,margin=2.2cm]{geometry}
\usetikzlibrary{arrows.meta,calc,angles,quotes,positioning,patterns,%
  backgrounds,shapes.geometric,decorations.pathmorphing,decorations.markings,intersections}
\definecolor{primary}{RGB}{222,49,99}\definecolor{primarydark}{RGB}{150,22,55}
\definecolor{defcol}{RGB}{0,114,178}\definecolor{methcol}{RGB}{52,92,148}
\definecolor{excol}{RGB}{0,135,90}\definecolor{attncol}{RGB}{200,40,55}
\tcbset{boxbase/.style={enhanced,breakable,boxrule=0.9pt,arc=2.5pt,
  left=3mm,right=3mm,top=2.3mm,bottom=2.3mm,fonttitle=\bfseries,coltitle=white}}
% --- boîtes TUTORIEL (noms EXACTS, ne pas renommer) ---
\newtcolorbox{cle}[1][]{boxbase,colback=defcol!6,colframe=defcol,
  title={\textsc{À retenir}\ifx\relax#1\relax\relax\else\textmd{ : \itshape #1}\fi}}
\newtcolorbox{methode}[1][]{boxbase,colback=methcol!7,colframe=methcol,sharp corners,
  title={\textsc{Méthode}\ifx\relax#1\relax\relax\else\textmd{ --- \itshape #1}\fi}}
\newtcolorbox{exemplebox}[1][]{boxbase,colback=excol!5,colframe=excol,
  title={\textsc{Exemple}\ifx\relax#1\relax\relax\else\textmd{ : \itshape #1}\fi}}
\newtcolorbox{piege}[1][]{boxbase,colback=attncol!6,colframe=attncol,
  title={\textsc{Piège}\ifx\relax#1\relax\relax\else\textmd{ : \itshape #1}\fi}}
\newtcolorbox{remarque}[1][]{boxbase,colback=black!4,colframe=black!45,
  title={\textsc{Remarque}\ifx\relax#1\relax\relax\else\textmd{ : \itshape #1}\fi}}
% --- boîtes EXERCICE / CORRIGÉ (noms EXACTS, auto-comptées) ---
\newtcolorbox[auto counter]{exo}[1][]{boxbase,colback=primary!4,colframe=primary,
  title={\textsc{Exercice}~\thetcbcounter\ifx\relax#1\relax\relax\else\textmd{ --- \itshape #1}\fi}}
\newtcolorbox[auto counter]{corrige}[1][]{boxbase,colback=excol!4,colframe=excol,
  title={\textsc{Corrigé}~\thetcbcounter\ifx\relax#1\relax\relax\else\textmd{ --- \itshape #1}\fi}}
\newcommand{\vv}[1]{\vec{#1}}\newcommand{\ds}{\displaystyle}
\newcommand{\R}{\mathbb{R}}\newcommand{\N}{\mathbb{N}}\newcommand{\Z}{\mathbb{Z}}
% (un \usepackage{fancyhdr}+en-tête est possible ; il est ignoré côté web)
% =========================================================================
\begin{document}
\sisetup{per-mode=symbol}

\begin{corrige}[compression maximale d'un ressort]
\begin{enumerate}[label=\alph*)]
  \item $v=\dfrac{36}{3,6}=\qty{10}{\metre\per\second}$, d'où
        $E_c=\tfrac12 m v^2=\tfrac12\times10\times10^2=\qty{500}{\joule}$.
  \item $\tfrac12 k x^2=E_c=500$, donc $x^2=\dfrac{2\times500}{k}=\dfrac{1000}{1000}=1$, soit
        $x=\qty{1}{\metre}$.
\end{enumerate}
\end{corrige}

\begin{corrige}[lecture d'un graphe force--élongation]
\begin{enumerate}[label=\alph*)]
  \item $x=\qty{10}{\centi\metre}=\qty{0,10}{\metre}$ et $F=\qty{2,5}{\newton}$, donc
        $k=\dfrac{F}{x}=\dfrac{2,5}{0,10}=\qty{25}{\newton\per\metre}$.
  \item $E_p=\tfrac12 k x^2=\tfrac12\times25\times(0,10)^2=\tfrac12\times25\times0,01=\qty{0,125}{\joule}$
        (c'est l'aire du triangle sous la droite).
  \item $F=kx$ \emph{croît} avec $x$ : la force à \qty{4}{\centi\metre} ($F=\qty{1}{\newton}$) est
        \textbf{inférieure} à celle à \qty{10}{\centi\metre} ($F=\qty{2,5}{\newton}$).
\end{enumerate}
\end{corrige}

\begin{corrige}[un ressort qui propulse un bloc]
\begin{enumerate}[label=\alph*)]
  \item $E_p=\tfrac12 k x^2=\tfrac12\times400\times(0,10)^2=\tfrac12\times400\times0,01=\qty{2}{\joule}$.
  \item $\tfrac12 m v^2=E_p=2\Rightarrow v^2=\dfrac{2\times2}{1}=4\Rightarrow v=\qty{2}{\metre\per\second}$.
\end{enumerate}
\end{corrige}

\begin{corrige}[vrai ou faux : énergies potentielles]
\begin{enumerate}[label=\alph*)]
  \item \textbf{Vrai}. $E_p=\tfrac12\times25\times(0,10)^2=\qty{0,125}{\joule}$.
  \item \textbf{Faux}. $F=kx$ augmente avec $x$ : $F(\qty{4}{\centi\metre})=\qty{1}{\newton}<F(\qty{10}{\centi\metre})=\qty{2,5}{\newton}$.
  \item \textbf{Faux}. $E_p\propto x^2$ : doubler $x$ \emph{quadruple} $E_p$ (facteur $2^2=4$).
  \item \textbf{Vrai}. $E_p=mgh$ dépend de la référence : si l'on place le « zéro » d'altitude
        \emph{au-dessus} de l'objet, sa hauteur $h$ est négative et $E_p<0$. Seule $\Delta E_p$ a un
        sens physique absolu.
\end{enumerate}
\end{corrige}
\end{document}
